
\begin{table}[H]
   \caption{\label{tab:neonatal_suggestive} \textbf{Early-term shifting and neonatal hospital use}}
   \centering
\small\setlength{\tabcolsep}{5pt}\resizebox{\ifdim\width>\linewidth \linewidth\else\width\fi}{!}{%
   \begin{tabular}{lccc}
      \toprule
 & (1)      & (2)      & (3)\\  
\addlinespace[2pt]
 & \multicolumn{2}{c}{Perinatal-condition admissions per private birth} & All infant admissions per private birth\\
      \midrule
      Share of private births at 37--38 weeks & 0.0044   &          & -0.0084\\   
                                              & (0.0097) &          & (0.0540)\\   
\addlinespace[2pt]
      For-profit cesarean rate                &          & 0.0094   &   \\   
                                              &          & (0.0070) &   \\   
\addlinespace[2pt]
      \midrule
      Municipality fixed effects & Yes      & Yes      & Yes\\  
      Year fixed effects & Yes      & Yes      & Yes\\  
      \midrule
      Observations                            & 4,284    & 4,284    & 4,284\\  
\bottomrule
   \end{tabular}
}
\begin{minipage}{\linewidth}\footnotesize
\textit{Notes:} Municipality-year cells, 2015--2024, weighted by private births; cells with at least 50 private births. Infant admissions are TISS hospital events of beneficiaries in the $<$1 age band, at the provider municipality; perinatal conditions are primary diagnoses of conditions originating in the perinatal period. Ecological and correlational, a corroboration of the early-term margin, not a causal estimate. Standard errors, clustered by municipality, are reported in parentheses. \footnotesize Significance levels: *** p$<$0.01, ** p$<$0.05, * p$<$0.10.
\end{minipage}
\end{table}


